%!TeX root=../main.tex

\chapter{Contesto e analisi del sistema}\label{ch:contesto-e-analisi-del-sistema}

\section{Il problema: frammentazione e inefficienza nella gestione dei contest}\label{sec:problema-frammentazione-inefficienza}
La gestione di un contest, specialmente in ambiti creativi, è un processo che coinvolge una molteplicità di attori (organizzatori, partecipanti, giurati) e fasi operative sequenziali e dipendenti.
L'approccio tradizionale si affida a un assemblaggio di strumenti generici e non integrati, creando un flusso di lavoro discontinuo, inefficiente e soggetto a errori.

Le criticità di questo modello possono essere analizzate nel dettaglio:
\begin{itemize}
    \item \textbf{Frammentazione del processo e dei dati}: la comunicazione degli inviti avviene tramite servizi di posta elettronica, la sottomissione dei lavori tramite posta elettronica o piattaforme di cloud storage (es.\ Google Drive) e la tabulazione dei voti tramite fogli di calcolo (es.\ Microsoft Excel).
    Questo approccio disperde le informazioni, rende difficile avere una visione d'insieme e richiede un costante lavoro manuale di organizzazione dei dati.
    \item \textbf{Mancanza di integrità e consistenza dei dati}: la gestione manuale aumenta esponenzialmente il rischio di errori.
    Esempi comuni includono la perdita di iscrizioni, l'errata trascrizione dei punteggi o l'inclusione di partecipanti non idonei.
    Non esiste un'unica ``fonte di verità'' per i dati del contest.
    \item \textbf{Carenza di strumenti in tempo reale}: l'assenza di un sistema centralizzato impedisce di gestire sessioni di voto live in modo coordinato e trasparente.
    \item \textbf{Rischi di sicurezza e potenziale di abuso}: la sottomissione di voti multipli, la partecipazione di utenti non invitati e la gestione di conflitti di interesse diventano problemi concreti e di difficile gestione.
\end{itemize}

L’obiettivo primario del sistema \textbf{Swift Contest} è superare queste criticità strutturali, fornendo un ambiente end-to-end integrato, sicuro e automatizzato, che funga da unica fonte di verità per l'intero ciclo di vita del contest.

\begin{figure}[H]
    \centering
    \includegraphics[width=\textwidth]{capitoli/figure/diagrams/traditional-vs-swiftcontest-flow}
    \caption{Confronto tra il flusso di lavoro tradizionale e l'approccio integrato di Swift Contest.}
    \label{fig:flow-comparison}
\end{figure}

\section{Stato dell’arte e analisi delle alternative}\label{sec:stato-dellarte-e-analisi-delle-alternative}
L'analisi dello stato dell'arte ha rivelato un'interessante lacuna nel panorama software attuale: \textbf{non esistono soluzioni verticali e accessibili che coprano in modo integrato l'intero ciclo di vita di un contest}.
Le pratiche correnti si affidano a due approcci principali, entrambi inadeguati.

\subsection{Strumenti generici adattati: lo ``Status Quo''}\label{subsec:strumenti-generici-adattati-lo-status-quo}
Questa è la pratica più diffusa e rappresenta il problema fondamentale che \textbf{Swift Contest} si propone di risolvere.
Consiste nell'utilizzare una combinazione di strumenti generici, non pensati per questo scopo, creando un flusso di lavoro frammentato e di difficile gestione, ad esempio:
\begin{itemize}
    \item \textbf{Google Forms} per la raccolta delle iscrizioni e la sottomissione dei voti.
    \item \textbf{Google Drive} per la sottomissione dei lavori.
    \item \textbf{Google Sheets} per la tabulazione manuale dei voti.
    \item \textbf{Email} per tutte le comunicazioni.
\end{itemize}
Come già analizzato, questo approccio soffre di gravi criticità in termini di sicurezza, integrità dei dati ed efficienza operativa, che hanno motivato la nascita del progetto.

\subsection{Piattaforme per domini correlati}\label{subsec:piattaforme-per-domini-correlati}
Sono state analizzate anche piattaforme software specializzate in domini vicini a quello dei contest, per verificare la possibilità di un loro adattamento.
Tuttavia, nessuna di esse si è rivelata una soluzione praticabile.

\subsubsection{Piattaforme di e-voting}
Sistemi come \textit{Helios Voting} sono progettati per il \textbf{voto politico o assembleare}.
Il loro focus è sulla crittografia e sull'anonimato del votante.
Di conseguenza, mancano completamente delle funzionalità necessarie per un contest, quali:
\begin{itemize}
    \item La possibilità di invitare partecipanti e creare giurie.
    \item La sottomissione e gestione di opere digitali.
    \item La creazione di form di voto personalizzati con campi testuali e numerici.
\end{itemize}
Il loro scopo è intrinsecamente diverso e non sono in alcun modo utilizzabili.

\subsubsection{Piattaforme per la gestione di eventi}
Soluzioni commerciali come \textit{Eventbrite} sono focalizzate sulla \textbf{gestione logistica di eventi e sulla vendita di biglietti}.
Sebbene eccellano nella gestione delle registrazioni, non offrono alcuna funzionalità specifica per il dominio dei contest, come:
\begin{itemize}
    \item La sottomissione di opere digitali.
    \item La creazione di giurie con ruoli definiti.
    \item La configurazione di form di valutazione strutturati.
    \item La gestione di sessioni di voto in tempo reale.
\end{itemize}
Il loro scopo è completamente diverso e non risultano realisticamente adattabili ai requisiti di un contest.

\subsection{Posizionamento di Swift Contest}\label{subsec:posizionamento-di-swift-contest}
L'analisi conferma che Swift Contest si posiziona in un \textbf{vuoto di mercato e tecnologico}.
Non mira a competere con le piattaforme di e-voting o di gestione eventi, ma a creare una \textbf{nuova categoria di prodotto}: una soluzione verticale, integrata e accessibile, progettata specificamente per le esigenze di organizzatori, partecipanti e giurati di un contest.

\section{Attori e casi d’uso}\label{sec:attori-e-casi-duso}
Il sistema è progettato per servire diversi tipi di utenti, ciascuno con un ruolo e un insieme di permessi ben definiti.
L'interazione tra questi attori e le funzionalità del sistema è stata modellata attraverso una serie di diagrammi dei Casi d'Uso, suddivisi in pacchetti logici (packages) che legano funzionalità della stessa categoria per migliorarne leggibilità e modularità.

\subsection{Attori del sistema}\label{subsec:attori-del-sistema}
L'analisi del dominio ha permesso di identificare i seguenti attori principali:
\begin{itemize}
    \item \textbf{Guest}: rappresenta un utente non autenticato.
    In questo sistema, che non prevede contenuti pubblici, un \emph{Guest} non ha privilegi di visualizzazione.
    Le sue uniche capacità sono quelle di \textbf{autenticarsi} (tramite registrazione o login) o di agire come \textbf{Simple Juror} se in possesso di un token di accesso.
    \item \textbf{Authenticated}: è un attore astratto che rappresenta un utente che ha completato l'autenticazione.
    Funge da base per i ruoli operativi, che ne ereditano le funzionalità comuni come la gestione del proprio profilo e delle notifiche.
    \item \textbf{Organizer}: è un utente autenticato responsabile della creazione e della gestione completa di un contest.
    \item \textbf{Participant}: è un utente autenticato che si iscrive a un contest per sottoporre il proprio lavoro.
    \item \textbf{Juror}: è un utente autenticato il cui ruolo primario è quello di far parte di una giuria nominata (\emph{appointed jury}), oppure votare come \emph{Simple Juror} se in possesso di un token per la votazione.
    \item \textbf{Admin}: è un ruolo di amministrazione, gestito tramite una dashboard esterna (Retool), con privilegi elevati per la manutenzione del sistema.
\end{itemize}

\subsection{Casi d'Uso: Gestione Account e Funzionalità Generali}\label{subsec:casi-uso-gestione-account-funzionalita}
Questa area copre le funzionalità di base legate all'accesso e alla gestione dell'account utente.
Include i processi di autenticazione, la modifica del profilo personale e la gestione delle notifiche interne.

\begin{figure}[H]
    \centering
    \includegraphics[width=0.6\textwidth]{capitoli/figure/diagrams/use-case-diagram-1}
    \caption{Diagramma dei Casi d'Uso: Gestione Account e Funzionalità Generali.}
    \label{fig:uc-account}
\end{figure}

\subsection{Casi d'Uso: Gestione del Ciclo di Vita del contest}\label{subsec:casi-uso-gestione-ciclo-vita-contest}
Questa è l'area funzionale più vasta e descrive il cuore del processo organizzativo.
Copre la creazione di un contest, la gestione dei suoi membri (partecipanti e giurati) e la sottomissione dei lavori.
Gli attori principali in questo contesto sono l'\emph{Organizer} e il \emph{Participant}.

\begin{figure}[H]
    \centering
    \includegraphics[width=0.6\textwidth]{capitoli/figure/diagrams/use-case-diagram-2}
    \caption{Diagramma dei Casi d'Uso: Gestione del Ciclo di Vita del contest.}
    \label{fig:uc-contest}
\end{figure}

\subsection{Casi d'Uso: Gestione delle Votazioni e dei Risultati}\label{subsec:casi-uso-gestione-votazioni-risultati}
Questa sezione si concentra interamente sul processo critico della votazione.
Include la configurazione del form di voto, la conduzione della sessione di voto (anche con restrizioni di geolocalizzazione) e la generazione e pubblicazione delle classifiche.
Gli attori principali qui sono l'\emph{Organizer} e l'\emph{Juror}.

\begin{figure}[H]
    \centering
    \includegraphics[width=0.5\textwidth]{capitoli/figure/diagrams/use-case-diagram-3}
    \caption{Diagramma dei Casi d'Uso: Gestione delle Votazioni e dei Risultati.}
    \label{fig:uc-voting}
\end{figure}

\subsection{Casi d'Uso: Amministrazione del Sistema}\label{subsec:casi-uso-amministrazione-sistema}
Quest'ultima area isola le funzionalità di amministrazione, che non sono accessibili tramite l'applicazione principale ma attraverso una dashboard esterna.
Queste operazioni sono riservate all'attore \emph{Admin} e sono fondamentali per la manutenzione e la moderazione della piattaforma.

\begin{figure}[H]
    \centering
    \includegraphics[width=0.4\textwidth]{capitoli/figure/diagrams/use-case-diagram-4}
    \caption{Diagramma dei Casi d'Uso: Amministrazione del Sistema.}
    \label{fig:uc-admin}
\end{figure}

\section{Requisiti del sistema}\label{sec:requisiti-del-sistema}
I requisiti del sistema sono stati classificati secondo il modello \textbf{FURPS+}, un'estensione del modello FURPS che include categorie aggiuntive per coprire in modo esaustivo tutti i vincoli di un progetto software.

\subsection{Requisiti funzionali (Functionality)}\label{subsec:requisiti-funzionali}
\begin{itemize}
    \item \textbf{RF-01: Gestione Account Utente}: Il sistema deve permettere a un utente di registrarsi (\emph{Sign Up}) e autenticarsi (\emph{Sign In}).
    \item \textbf{RF-02: Gestione del Profilo Personale}: Ogni utente autenticato deve poter visualizzare e modificare i dati del proprio profilo (es.\ nome completo).
    Deve inoltre essere garantita la possibilità di eliminare il proprio account in modo autonomo (\emph{self-service deletion}).
    \item \textbf{RF-03: Gestione Contest}: Il sistema deve permettere a un \emph{Organizer} di creare, modificare e cancellare i propri contest.
    \item \textbf{RF-04: Gestione Giurie e Form di Voto}: Il sistema deve permettere a un \emph{Organizer} di creare giurie (di tipo \emph{appointed} o \emph{simple}) e di configurare un \emph{VotingForm} associato, con campi organizzati secondo il loro \textbf{scope} e tipo.
    \item \textbf{RF-05: Gestione Iscrizioni}: Il sistema deve permettere a un \emph{Organizer} di invitare partecipanti e giurati nominati tramite email con token univoci.
    \item \textbf{RF-06: Gestione Esclusioni di Voto}: Il sistema deve permettere a un \emph{Organizer} di configurare delle esclusioni, impedendo a specifici \emph{Appointed Juror} di valutare determinati \emph{Participant} all'interno di una sessione di voto.
    \item \textbf{RF-07: Accettazione Inviti e Accesso al contest}: Il sistema deve permettere a un utente di utilizzare un token di invito per formalizzare la propria iscrizione a un contest, sia come \emph{Participant} che come \emph{Appointed Juror}.
    \item \textbf{RF-08: Sottomissione Lavori}: Il sistema deve permettere a un \emph{Participant} di sottomettere un proprio \emph{Work} (Lavoro Digitale) per un contest a cui è iscritto.
    \item \textbf{RF-09: Restrizione Geografica della Votazione}: Il sistema deve permettere a un \emph{Organizer} di imporre una restrizione geografica a una sessione di voto, richiedendo ai giurati di trovarsi entro un raggio specifico da un luogo predefinito.
    \item \textbf{RF-10: Processo di Votazione}: Il sistema deve permettere a un utente autorizzato (sia \emph{Appointed Juror} che \emph{Simple Juror}) di inviare i propri voti durante una sessione di voto attiva, rispettando le eventuali esclusioni configurate.
    \item \textbf{RF-11: Gestione Classifiche e Risultati}: Il sistema deve permettere all'\emph{Organizer} di generare una classifica per ogni singola giuria, di esportare i risultati grezzi in formato Excel e di pubblicare una classifica finale in formato PDF\@.
    \item \textbf{RF-12: Visualizzazione dei Risultati}: Il sistema deve permettere a \emph{Participant} e \emph{Juror} di un contest di visualizzare e scaricare le classifiche finali pubblicate dall'organizzatore.
    \item \textbf{RF-13: Notifiche Inbox}: Il sistema deve fornire a ogni utente autenticato una inbox per ricevere notifiche generate automaticamente da eventi di sistema.
    \item \textbf{RF-14: Amministrazione del Sistema}: Il sistema deve esporre un'interfaccia sicura per permettere a un \emph{Admin} di eseguire operazioni di alto livello, come la visualizzazione di tutti gli utenti o la rimozione forzata di un contest.
\end{itemize}

\subsection{Requisiti non funzionali}\label{subsec:requisiti-non-funzionali}
\begin{itemize}
    \item \textbf{Usabilità (Usability)}: L'interfaccia utente deve essere intuitiva, facile da apprendere e coerente tra la versione mobile e quella web.
    \item \textbf{Affidabilità (Reliability)}: Il sistema deve essere robusto e garantire l'integrità dei dati.
    La logica di \emph{snapshot} delle sessioni di voto è un requisito chiave per assicurare che il processo di valutazione sia immune da modifiche ai dati.
    \item \textbf{Performance}: L'applicazione deve essere reattiva, con tempi di caricamento inferiori ai 3 secondi.
    %L'uso mirato del servizio Realtime è una scelta progettuale per minimizzare il traffico di rete.
    \item \textbf{Manutenibilità (Supportability)}: L'architettura software deve essere modulare.
    L'adozione del pattern MVVM/BLoC per il frontend garantisce una netta separazione delle responsabilità, facilitando futuri interventi.
    \item \textbf{Implementazione (Implementation)}:
    \begin{itemize}
        \item Il client deve essere sviluppato in \textbf{Dart} con il framework \textbf{Flutter} per le piattaforme Android e Web.
        \item Il backend deve basarsi sulla piattaforma \textbf{Supabase}, utilizzando \textbf{PostgreSQL}, \textbf{PL/pgSQL} per le RPC e \textbf{TypeScript} per le Edge Functions.
        \item La gestione dello stato nel client deve seguire il pattern \textbf{BLoC}.
    \end{itemize}
    \item \textbf{Interfaccia (Interface)}:
    \begin{itemize}
        \item Il sistema deve interfacciarsi in modo sicuro con servizi esterni come \textbf{Resend} (per le email) e \textbf{Google Places} (per la geolocalizzazione).
        \item Deve essere prevista un'API sicura per l'interfacciamento con strumenti di amministrazione esterni (\textbf{Retool}).
    \end{itemize}
    \item \textbf{Operativi (Operation)}:
    \begin{itemize}
        \item Il deployment del backend deve essere tracciabile tramite migrazioni SQL gestite dalla CLI di Supabase.
        \item Il deployment del frontend web deve essere automatizzato tramite una pipeline di CI/CD\@.
        \item I segreti (API keys) devono essere gestiti in modo sicuro, senza essere esposti nel codice sorgente.
    \end{itemize}
    \item \textbf{Distribuzione (Packaging)}:
    \begin{itemize}
        \item L'applicazione web deve essere accessibile tramite URL pubblico.
        \item L'applicazione Android deve essere distribuita come pacchetto \textbf{APK}, tramite un meccanismo di download sicuro.
    \end{itemize}
    \item \textbf{Legali (Legal)}:
    \begin{itemize}
        \item Il sistema deve essere sviluppato utilizzando principalmente tecnologie con licenze open-source permissive.
        \item La gestione dei dati personali deve essere orientata ai principi del GDPR, che garantisce all'utente il diritto alla cancellazione.
    \end{itemize}
\end{itemize}

\section{Modello del dominio}\label{sec:modello-del-dominio}
Il dominio applicativo è stato modellato a partire dalle entità fondamentali e dalle loro relazioni.
Il Diagramma di Classe Concettuale (UML) è astratto dai dettagli implementativi e si focalizza sui concetti di business e sui vincoli che li governano.

\begin{figure}[H]
    \centering
    \includegraphics[width=0.8\textwidth]{capitoli/figure/diagrams/class-diagram}
    \caption{Diagramma delle classi}
    \label{fig:class-diagram}
\end{figure}

\subsection{Entità fondamentali}\label{subsec:entita-fondamentali}
\begin{itemize}
    \item \textbf{User}: l'utente del sistema, contenente informazioni del profilo.
    \item \textbf{Contest}: l'entità centrale, che aggrega tutte le informazioni relative a un concorso.
    \item \textbf{ParticipantMembership} e \textbf{JurorMembership}: classi associative che risolvono le relazioni ``molti-a-molti'' tra utenti e contest/giurie.
    \item \textbf{Work}: l'elaborato unico caricato da un partecipante.
    È legato a una \emph{ParticipantMembership}.
    \item \textbf{Jury}: una giuria all'interno di un contest, che può essere di tipo \emph{appointed} o \emph{simple}.
    \item \textbf{VotingForm} e \textbf{VotingFormField}: la struttura del form di votazione.
    L'attributo \emph{scope} del campo ne determina il comportamento.
    \item \textbf{VotingSession}: una sessione di voto live.
    Al suo avvio, il sistema ``congela'' lo stato del contest creando delle righe dedicate in tabelle di sessione (es.\ \texttt{voting\_session\_participants}), che contengono non solo i riferimenti alle entità originali, ma anche una copia di \textbf{attributi di snapshot} (es.\ il nome del partecipante, il nome dell'opera) per garantire l'immutabilità del processo.
    \item \textbf{ContestRanking}: la classifica finale, rappresentata come un riferimento a un file nello storage.
    \item \textbf{Message}: una notifica di sistema inviata a un utente.
\end{itemize}

\subsection{Vincoli chiave del dominio}\label{subsec:vincoli-chiave-del-dominio}
Il modello dati è stato progettato per far rispettare programmaticamente alcune regole di business fondamentali:
\begin{enumerate}
    \item \textbf{Unicità del Work}: un partecipante può sottomettere un solo \emph{Work} per ogni concorso a cui è iscritto, garantito da un vincolo di unicità a livello di database.
    \item \textbf{Calcolo della Classifica Semplice}: la classifica viene generata esclusivamente sulla base dei campi di tipo \emph{slider} con scope \emph{participant}.
    I campi testuali non contribuiscono al punteggio.
    \item \textbf{Immutabilità della Votazione}: il meccanismo di snapshot garantisce che una volta avviata una sessione di voto, i dati rilevanti dei partecipanti, dei giurati e la struttura del form vengano ``congelati''.
    Questo avviene tramite la copia di attributi chiave in tabelle di sessione dedicate, rendendo il processo di valutazione immune da modifiche successive ai dati originali.
    \item \textbf{Accesso Ospite Controllato}: l'accesso come \emph{Simple Juror} è strettamente controllato da un token associato a una specifica giuria di tipo semplice (simple jury).
\end{enumerate}
